\documentclass{article}
\usepackage[T1]{fontenc}
\usepackage{lmodern}
\usepackage{graphicx}
\usepackage{amsmath,amssymb}
\usepackage[a4paper,margin=2cm]{geometry}
\usepackage[absolute,overlay]{textpos}
\setlength{\TPHorizModule}{1mm}
\setlength{\TPVertModule}{1mm}
\usepackage[indonesian]{babel}
\usepackage{hyperref}
\setlength{\emergencystretch}{2em}
% unit: Halaman 3

\begin{document}

% === Halaman 3 (llm) ===

\begin{itemize}
    \item Meskipun disebut kriptografi modern, namun algoritmanya tetap menggunakan dua teknik dasar dasar di dalam kriptografi klasik: \textbf{teknik substitusi} dan \textbf{teknik transposisi (permutasi)},
    \item tetapi operasinya dibuat lebih kompleks, tidak sesederhana \textit{cipher} klasik. Tujuannya: agar \textit{cipher} modern lebih sulit dikriptanalisis
    \item Selain kedua teknik dasar tersebut, juga digunakan teknik lain seperti rotasi, kompresi, ekspansi, penjumlahan modulo, dan lain-lain.
    \item Kriptografi modern melahirkan konsep-konsep baru seperti algoritma kriptografi kunci-publik, fungsi \textit{hash}, protokol kriptografi, tanda-tangan digital, pembangkit bilangan acak, skema pembagian kunci, dsb.
\end{itemize}

\end{document}
